\section*{CHAPTER 2: REVIEW OF RELATED LITERATURE AND STUDIES}
\addcontentsline{toc}{section}{CHAPTER 2: REVIEW OF RELATED LITERATURE AND STUDIES}

\subsection*{2.1 The Digital Divide and Infrastructure Deficit}
The ``digital divide'' in the Philippines is often a matter of hardware access rather than just connectivity. Technical learning (e.g., programming, CAD, graphics) requires desktop computers, which public institutions struggle to procure and maintain due to rapid depreciation and high CapEx requirements \cite{ulas2018}.

\subsection*{2.2 The Economic Decline of the Internet Café Industry}
Market analyses indicate a severe decline in the internet café sector as users shift to mobile platforms \cite{garcia2019}. While these shops possess high-end hardware for esports, they face heavy downtime during the day, making them prime candidates for government resource-sharing models.

\subsection*{2.3 Cybercafes as Informal Learning Spaces (ICT4D)}
Literature in ICT for Development (ICT4D) recognizes cybercafes as vital community hubs \cite{robinson2009}. However, they are currently used as informal student workarounds. This study proposes moving from informal usage to institutionalized governance.

\subsection*{2.4 The Philippine PPP Code (RA 11966)}
The newly enacted PPP Code of 2023 \cite{pppcode2023} recognizes the autonomy of LGUs to enter into local PPP projects. It emphasizes Value for Money (VFM), allowing LGUs to lease private assets to deliver public services (digital education) more efficiently than traditional procurement.

\subsection*{2.5 TESDA Modalities and Precedents}
TESDA’s Community-Based Training (CBT) and Educational Service Contracting (ESC) provide precedents for the government paying private partners to deliver public education. This study applies this ``infrastructure service contracting'' logic to the ICT sector.